% vim: ai sw=2 spell spelllang=cs

\section{Objektové soubory}
\frame{\sectionpage}

\begin{frame}
  \frametitle{Objektové soubory}

  \begin{itemize}
    \item Binární soubory s kódem (a daty)
    \item Mohou mít různý formát

    \pause

    \item Výstup překladače
    \begin{itemize}
      \item \cmd{gcc -c} produkuje \file{.o}
      \item Obvykle nespustitelný objekt
    \end{itemize}

    \pause

    \item Výstup linkeru
    \begin{itemize}
      \item Statická/dynamická knihovna
      \item Spustitelná binárka
    \end{itemize}
  \end{itemize}

  \begin{textblock*}{4cm}(\textwidth-6cm,2.8cm)
    \only<2->{
    \tikzstyle{pre} = [<-,shorten <=1pt,>=stealth',semithick]
    \begin{tikzpicture}[node distance=5mm, font=\footnotesize]
	\path[use as bounding box] (-2,0) rectangle (1,-2);
	\node [tl_phase] (GCC1) {\cmd{gcc}};
	\node<3-> [tl_phase] (LD1) [left=of GCC1] {\cmd{ld}};
	\node [tl_phase] (GCC2) [right=of GCC1] {\cmd{gcc}};

	\node<3-> [tl_state] (SO1) [below=of LD1] {\file{.so}}
	  edge [pre] node {} (LD1);
	\node [tl_state] (OBJ1) [below=of GCC1] {\file{.o}}
	  edge [pre] node {} (GCC1);
	\node [tl_state] (OBJ2) [below=of GCC2] {\file{.o}}
	  edge [pre] node {} (GCC2);

	\node<3-> [tl_phase] (LD2) [below=of OBJ1] {\cmd{ld}}
	  edge [pre] node {} (SO1)
	  edge [pre] node {} (OBJ1)
	  edge [pre] node {} (OBJ2);

	\node<3-> [tl_state] (A) [below=3cm of LD1] {\file{.a}}
	  edge [pre] node {} (LD2);
	\node<3-> [tl_state] (SO2) [below=3cm of GCC1] {\file{.so}}
	  edge [pre] node {} (LD2);
	\node<3-> [tl_state] (BIN) [below=3cm of GCC2] {\file{bin}}
	  edge [pre] node {} (LD2);
    \end{tikzpicture}
  }
  \end{textblock*}

\end{frame}

\begin{frame}
  \frametitle{Obsah objektových souborů}

  \begin{itemize}
    \item Obvykle
    \begin{itemize}
      \item Identifikace souboru
      \begin{itemize}
	\item ELF, PE, \dots
      \end{itemize}

      \pause

      \item Informace o souboru
      \begin{itemize}
	\item Cílová architektura, struktura, \dots
      \end{itemize}

      \pause

      \item Kód
      \item Data

      \pause

      \item Ostatní informace potřebné k běhu/linkování
      \begin{itemize}
	\item Relokace, informace o symbolech atd.
      \end{itemize}
    \end{itemize}

    \pause

    \item Ale také např.
    \begin{itemize}
      \item Ladicí informace
      \item Informace o překladači
    \end{itemize}

    \pause

    \item Tyto informace bývají v oddílech -- tzv.\ sekcích

  \end{itemize}

\end{frame}

\begin{frame}
  \frametitle{Úkol}

  \heading{Práce s objektovým souborem}
  \begin{enumerate}
    \item Vyberte si nějaký objektový soubor
    \begin{itemize}
      \item Např.\ \file{/bin/true}
    \end{itemize}
    \item Vypište si informace o souboru
    \begin{itemize}
      \item \cmd{objdump -f}
    \end{itemize}
    \item Vypište si seznam sekcí
    \begin{itemize}
      \item \cmd{objdump -h}
    \end{itemize}
    \item Vypište si dump sekce \code{.rodata}
    \begin{itemize}
      \item \cmd{objdump -s -j .rodata}
    \end{itemize}
    \item Vypište si debug info
    \begin{itemize}
      \item \cmd{objdump -Wi}
      \item Pravděpodobně se nevypíše nic -- proč?
    \end{itemize}
  \end{enumerate}

\end{frame}

\begin{frame}
  \frametitle{Formáty objektových souborů}

  \begin{itemize}
    \item Několik forem/standardů
    \begin{itemize}
      \item a.out (UNIX/starý Linux)
      \item \emph{ELF} (Linux)
      \item COM (DOS, čistě kód+data, nic víc)
      \item COFF (UNIX, ELF ho převálcoval)
      \item MZ (DOS .exe)
      \item PE (Windows .exe, vychází z COFF)
    \end{itemize}
  \end{itemize}

  \pause

  \heading{Demo:} \textsc{DosBox} a reboot přes COM (\code{ea f0 ff 00 f0})

\end{frame}

\begin{frame}
  \frametitle{Úkol}

  \heading{Objdump neumí jen ELF}
  \begin{enumerate}
    \item Stáhněte si nějakou Windows aplikaci
    \begin{itemize}
      \item Cokoliv, co je .exe (PE), např. putty
    \end{itemize}
    \item Vypište si informace o souboru
    \begin{itemize}
      \item \cmd{objdump -f}
    \end{itemize}
    \item Vypište si informace závislé na formátu
    \begin{itemize}
      \item \cmd{objdump -p}
      \item Zjistěte DLL, na kterých váš program závisí
    \end{itemize}
    \item Porovnejte sekce s předchozím úkolem
  \end{enumerate}

\end{frame}

\section{\file{libbfd}}
\frame{\sectionpage}

\begin{frame}
  \frametitle{\file{libbfd}}

  \begin{itemize}
    \item Knihovna pro práci s \emph{různými} binárními formáty
      (\header{bfd.h}, v binutils)
    \begin{itemize}
      \item a.out, ELF, PE, binární, \dots
      \item Seznam: \code{const char **bfd_target_list()}
    \end{itemize}

    \pause

    \item A různými architekturami
    \begin{itemize}
      \item Seznam: \code{const char **bfd_arch_list()}
    \end{itemize}

    \pause

    \item Dokumentace
    \begin{itemize}
      \item \url{https://sourceware.org/binutils/docs/bfd/}
      \item \emph{Některá užitečná makra bez dokumentace} (jen v \header{bfd.h})
    \end{itemize}

    \pause

    \item Knihovny při překladu: \cmd{gcc ... -lbfd -liberty -ldl -lz}

    \pause

    \item Nutno inicializovat pomocí \code{void bfd_init()}
  \end{itemize}

\end{frame}

\begin{frame}[fragile]{Chyby}

  \heading{Jak zjistit, co se stalo?}
  \begin{itemize}
    \item Poslední chyba: \code{bfd_error_type bfd_get_error()}
    \item Převod na text: \code{const char *bfd_errmsg(bfd_error_type error_tag)}
  \end{itemize}

  \bigskip

  \begin{block}{Typicky}
  \begin{lstlisting}
if (!bfd_function(...))
  errx(1, "bfd_function: %s", bfd_errmsg(bfd_get_error()));
  \end{lstlisting} %stopzone
  \end{block}

\end{frame}

\begin{frame}
  \frametitle{Úkol}

  \heading{Výpis podporovaných formátů a architektur vaší \file{libbfd}}
  \begin{enumerate}
    \item Doplňujte \file{pb173-bin/03/bfd.c}
    \item Zavolejte \code{bfd_target_list}
    \item Vypište návratové hodnoty (v cyklu)
    \begin{itemize}
      \item Dokud nenarazíte na \code{NULL}
    \end{itemize}
    \item Zavolejte \code{bfd_arch_list}
    \item Vypište návratové hodnoty (v cyklu)
    \begin{itemize}
      \item Dokud nenarazíte na \code{NULL}
    \end{itemize}
    \item Přeložte a spusťte
  \end{enumerate}

\end{frame}

\begin{frame}[fragile]{Soubory}

  \begin{itemize}
    \item Otevření
    \begin{itemize}
      \item \code{bfd *bfd_openr(const char *file, const char *target)}
      \item \code{bfd *bfd_openw(const char *file, const char *target)}
      \item Jako \code{target} používejte \code{NULL}
      \item Návratová hodnota se používá jako držátko
    \end{itemize}

    \pause

    \item Ověření
    \begin{itemize}
      \item \emph{Nutné před prací se souborem}
      \item \code{bfd_boolean bfd_check_format(bfd *b, bfd_format fmt)}
      \item Jako \code{fmt} používejte \code{bfd_object}
    \end{itemize}

    \pause

    \item Zavření
    \begin{itemize}
      \item \code{bfd_boolean bfd_close(bfd *b)}
    \end{itemize}
  \end{itemize}

  \pause

  \begin{block}{Příklad}
  \begin{lstlisting}[aboveskip=2pt,belowskip=2pt]
b = bfd_openr("file", NULL);
if (!bfd_check_format(b, bfd_object))
  warnx("Not an object!");
bfd_close(b);
  \end{lstlisting}
  \end{block}

\end{frame}

\begin{frame}{Úkol}

  \heading{Doplňte otevírání souboru zadaného jako parametr programu}
  \begin{enumerate}
    \item Zavolejte \code{bfd_openr}
    \item Zavolejte \code{bfd_check_format}
    \item Zavolejte \code{bfd_close}

    \item Ověřujte návratové hodnoty podle manuálu
    \begin{itemize}
      \item Vypisujte chyby při neúspěchu
    \end{itemize}

    \item Přeložte a spusťte
  \end{enumerate}

\end{frame}

\begin{frame}{Sekce}

  \begin{itemize}
    \item Pro každou existuje struktura \code{asection}
    \begin{itemize}
      \item Jméno: \code{bfd_get_section_name}
      \item Velikost: \code{bfd_get/set_section_size}
      \item Obsah: \code{bfd_get/set_section_contents}
      \item A další: \code{flags}, \code{vma}, \code{lma}, \dots
    \end{itemize}

    \pause

    \item Nalezení: \code{bfd_get_section_by_name}

    \pause

    \item Iterace: \code{bfd_map_over_sections}
    \begin{itemize}
      \item Háček, který se zavolá pro (skoro) každou sekci
    \end{itemize}

    \pause

    \item Vytvoření: \code{bfd_make_section_with_flags}
    \begin{itemize}
      \item Flags: \code{SEC_HAS_CONTENTS}, pokud má sekce i obsah
    \end{itemize}
  \end{itemize}

\end{frame}

\begin{frame}{Úkol}

  \heading{Hexdump sekcí (součást domácího)}
  \begin{enumerate}
    \item Iterujte přes sekce (\code{bfd_map_over_sections})
    \item Vypište údaje o každé sekci
    \begin{itemize}
      \item Název
      \item Velikost
      \item Flags
      \item Hexdump prvních 16 bytů obsahu
    \end{itemize}
    \item Přeložte
    \item Spusťte
    \begin{itemize}
      \item Pro ELF
      \item Pro PE
    \end{itemize}
  \end{enumerate}

\end{frame}
